\section{Data Recipe：数据配方才是壁垒}

同样的数据源，不同的清洗、过滤、配比、训练顺序和反馈方式，会得到不同模型能力。Recipe 是数据产业中最隐性的壁垒。它不是“买来一堆数据”，而是决定哪些数据进训练、以什么比例进、何时进、如何评估。

\begin{figure}[H]
\centering
\includegraphics[width=0.92\textwidth]{data-recipe.png}
\caption{Data Recipe：来源、过滤、配比、课程、评估反馈。}
\end{figure}

\begin{knowledgebox}{读图：Recipe 为什么比原料更难复制}
图中中心是 Data Recipe，周围是来源、过滤、配比、课程和评估。原料可以购买，配方来自实验、失败、评估和领域知识。Recipe 决定数据是否真正转化成能力。
\end{knowledgebox}

\subsection{术语消化：Recipe 相关概念}

\noindent\begin{tabularx}{\textwidth}{p{0.22\textwidth}p{0.32\textwidth}X}
\toprule
术语 & 解决的问题 & 本期中的含义 \\
\midrule
Filtering & 去掉低质、重复、有害或无关数据 & 决定训练信号纯度。 \\
Mixture & 不同来源数据的比例 & 决定能力分布和偏好。 \\
Curriculum & 数据进入训练的顺序 & 决定模型先学什么、后学什么。 \\
Eval Feedback & 用评测反推数据改动 & 让 recipe 可迭代。 \\
Synthetic Data & 由模型或仿真生成的数据 & 解决真实数据稀缺，但需控制分布偏差。 \\
\bottomrule
\end{tabularx}

\subsection{本章小结}

Data Recipe 是从数据到能力的转换器。数据产业真正难复制的部分，常常不是数据源本身，而是配方和反馈闭环。

